\section{Thiết kế dữ liệu}

Mô hình dữ liệu được thiết kế để lưu cấu hình Feature Flag, phục vụ Evaluation API và ghi nhận lịch sử thay đổi. Nội dung phần này chỉ mô tả cấu trúc lưu trữ; thứ tự đánh giá nằm ở mục Rule Evaluation Engine và contract HTTP nằm ở mục API.

\subsection{Mô hình dữ liệu tổng thể}

Các thực thể chính của hệ thống được tóm tắt trong Bảng~\ref{tab:quan-he-thuc-the}. PostgreSQL được lựa chọn vì hỗ trợ giao dịch, khóa ngoại và kiểu \texttt{JSONB} cho rule parameters hoặc audit snapshot \cite{postgresql_json_types}.

\begin{table}[H]
    \centering
    \caption{Các thực thể và quan hệ chính trong mô hình dữ liệu}
    \label{tab:quan-he-thuc-the}
    \begin{tabular}{p{0.26\textwidth} p{0.64\textwidth}}
        \toprule
        \textbf{Thực thể} & \textbf{Vai trò chính} \\
        \midrule
        \texttt{Project} & Phạm vi quản lý cao nhất của một ứng dụng/sản phẩm. \\
        \texttt{Environment} & Môi trường triển khai như staging hoặc production. \\
        \texttt{FeatureFlag} & Định nghĩa ổn định của flag trong project. \\
        \texttt{FlagEnvironmentConfig} & Cấu hình runtime của một flag trong một environment. \\
        \texttt{FlagRule} & Rule targeting thuộc một cấu hình runtime. \\
        \texttt{FlagGroup} & Nhóm các flag liên quan trong cùng project. \\
        \texttt{FlagGroupConfig} & Trạng thái group kill switch theo environment. \\
        \texttt{SampleUserContext} & Dữ liệu người dùng mẫu phục vụ demo/kiểm thử. \\
        \texttt{AuditLogEntry} & Bản ghi append-only cho mutation quản trị. \\
        \bottomrule
    \end{tabular}
\end{table}

\subsubsection{Các thực thể chính}

Các thực thể được tổ chức theo chuỗi: project chứa environment và flag; flag có cấu hình theo environment; cấu hình có rule; group và audit log bổ sung khả năng vận hành/truy vết.

\subsubsection{Quan hệ giữa các thực thể}

Quan hệ quan trọng nhất là \texttt{FeatureFlag} -- \texttt{FlagEnvironmentConfig} -- \texttt{Environment}, cho phép một flag có cấu hình khác nhau theo từng môi trường. Group và audit log đều nằm trong phạm vi project để tránh lẫn dữ liệu giữa các ứng dụng.

\subsubsection{Phân tách định nghĩa flag và cấu hình runtime}

\texttt{FeatureFlag} lưu danh tính, tên, mô tả và lifecycle; \texttt{FlagEnvironmentConfig} lưu enabled/disabled, serving mode và kill switch theo environment. Việc tách này tránh nhân bản flag khi production và staging cần cấu hình khác nhau.

\subsubsection{Sơ đồ thực thể liên kết}

Sơ đồ rút gọn trong Hình~\ref{fig:feature-flag-er-diagram} nhấn mạnh đường đi dữ liệu phục vụ evaluation.

\begin{figure}[H]
    \centering
    \begin{tikzpicture}[
        entity/.style={rectangle, draw, rounded corners, minimum width=3.1cm, minimum height=0.9cm, text width=3.0cm, align=center},
        relation/.style={-{Latex[length=2mm]}, thick},
        node distance=1.4cm and 1.7cm
    ]
        \node[entity] (project) {\textbf{Project}};
        \node[entity, below left=of project] (environment) {\textbf{Environment}};
        \node[entity, below right=of project] (flag) {\textbf{FeatureFlag}};
        \node[entity, below=2.4cm of project] (config) {\textbf{FlagEnvironmentConfig}};
        \node[entity, below=of config] (rule) {\textbf{FlagRule}};
        \node[entity, left=of config] (group) {\textbf{FlagGroup / Config}};
        \node[entity, right=of config] (audit) {\textbf{AuditLogEntry}};
        \draw[relation] (project) -- node[left] {\small 1:N} (environment);
        \draw[relation] (project) -- node[right] {\small 1:N} (flag);
        \draw[relation] (environment) -- node[left] {\small 1:N} (config);
        \draw[relation] (flag) -- node[right] {\small 1:N} (config);
        \draw[relation] (config) -- node[right] {\small 1:N} (rule);
        \draw[relation] (project) -- node[above] {\small 1:N} (group);
        \draw[relation] (project) -- node[above] {\small 1:N} (audit);
    \end{tikzpicture}
    \caption{Sơ đồ thực thể liên kết rút gọn của hệ thống Feature Flag}
    \label{fig:feature-flag-er-diagram}
\end{figure}

\subsection{Project và environment}

Project xác định phạm vi nghiệp vụ; environment xác định phạm vi triển khai. Khóa nghiệp vụ của hai thực thể này được dùng trong API, cache key và seed data.

\subsubsection{Thực thể Project}

\texttt{Project} gồm \texttt{id}, \texttt{key}, \texttt{name}, \texttt{description}, \texttt{createdAt} và \texttt{updatedAt}. \texttt{key} là duy nhất toàn hệ thống và ổn định trong vòng đời project.

\subsubsection{Thực thể Environment}

\texttt{Environment} gồm \texttt{id}, \texttt{projectId}, \texttt{key}, \texttt{name}, \texttt{description}, \texttt{isDefault} và timestamp. \texttt{key} là duy nhất trong phạm vi project.

\subsubsection{Environment mặc định}

Mỗi project có một environment mặc định để tạo cấu hình ban đầu cho flag mới và giúp demo chạy ngay sau seed.

\subsubsection{Ràng buộc environment trong project}

Environment luôn thuộc một project cụ thể; các truy vấn evaluation dùng cặp \texttt{projectKey}/\texttt{environmentKey} để tránh nhầm cấu hình giữa project.

\subsection{Feature Flag và cấu hình theo environment}

Thiết kế tách metadata của flag khỏi trạng thái runtime nhằm hỗ trợ nhiều environment.

\subsubsection{Thực thể FeatureFlag}

\texttt{FeatureFlag} gồm \texttt{id}, \texttt{projectId}, \texttt{key}, \texttt{name}, \texttt{description}, \texttt{status} và timestamp.

\subsubsection{Trạng thái vòng đời của flag}

Lifecycle status gồm active và archived. Archived giữ lịch sử nhưng không còn phục vụ runtime.

\subsubsection{Thực thể FlagEnvironmentConfig}

\texttt{FlagEnvironmentConfig} nối flag với environment và lưu \texttt{enabled}, \texttt{servingMode}, \texttt{killSwitch} cùng timestamp.

\subsubsection{Trạng thái cấu hình và serving mode}

Config disabled dừng flag trong environment; serving mode quyết định bật toàn cục hay xét rule khi config enabled.

\subsubsection{Cấu hình flag theo từng environment}

Cùng một flag có thể bật ở staging, rollout một phần ở production và disabled ở environment khác mà không cần tạo nhiều flag.

\subsection{Rule và tham số rule}

Rule được lưu riêng để có thể thay đổi chiến lược rollout mà không sửa định nghĩa flag.

\subsubsection{Thực thể FlagRule}

\texttt{FlagRule} gồm \texttt{id}, \texttt{flagConfigId}, \texttt{type}, \texttt{priority}, \texttt{enabled}, \texttt{parameters} và timestamp.

\subsubsection{Kiểu rule}

Mini project hỗ trợ ba kiểu: \texttt{USER\_ALLOWLIST}, \texttt{ROLE\_TARGETING} và \texttt{PERCENTAGE\_ROLLOUT}.

\subsubsection{Thứ tự ưu tiên của rule}

Database lưu \texttt{priority} cho rule cùng loại; type precedence của engine được định nghĩa riêng trong mục evaluation.

\subsubsection{Lưu trữ tham số bằng JSON}

\texttt{parameters} dùng \texttt{JSONB} để lưu danh sách user, danh sách role hoặc tỷ lệ rollout mà không cần tạo bảng phụ cho từng kiểu rule.

\subsubsection{Validation tham số theo loại rule}

Backend kiểm tra shape của \texttt{parameters}: user/role phải là danh sách chuỗi, percentage phải nằm trong miền hợp lệ.

\subsection{User context mẫu}

\texttt{SampleUserContext} phục vụ demo và kiểm thử, không phải nguồn danh tính bắt buộc của Evaluation API.

\subsubsection{Targeting key}

\texttt{targetingKey} là khóa ổn định cho percentage rollout và là duy nhất trong project demo.

\subsubsection{Các trường nghiệp vụ và attributes}

Dữ liệu mẫu có thể chứa \texttt{userId}, \texttt{role} và \texttt{attributes} như organization ID.

\subsubsection{Phạm vi sử dụng dữ liệu mẫu}

Dữ liệu này giúp chọn nhanh context trên Dashboard/Demo Application và không đại diện cho user store sản xuất.

\subsection{Flag group và group kill switch}

Flag group hỗ trợ quản lý các flag liên quan và cho phép ngắt nhanh cả nhóm trong một environment.

\subsubsection{Thực thể FlagGroup và FlagGroupConfig}

\texttt{FlagGroup} lưu định danh nhóm; \texttt{FlagGroupConfig} lưu kill switch của nhóm theo environment.

\subsubsection{Quan hệ giữa group và flag}

Một group có nhiều flag; trong phạm vi mini project, một flag thuộc tối đa một group.

\subsubsection{Trạng thái group kill switch theo environment}

Kill switch của group được lưu theo environment để tránh ảnh hưởng đồng thời staging và production.

\subsubsection{Audit log và cache invalidation cho group}

Mutation group tạo audit entry và invalidation snapshot của các flag thành viên bị ảnh hưởng.

\subsection{Audit log và lịch sử cấu hình}

Audit log là nguồn truy vết chính cho mọi mutation quản trị và cũng là cơ sở cho audit-backed history \cite{owasp_logging}.

\subsubsection{Cấu trúc AuditLogEntry}

\texttt{AuditLogEntry} gồm \texttt{id}, \texttt{projectId}, \texttt{actor}, \texttt{requestId}, \texttt{targetType}, \texttt{targetId}, \texttt{targetKey}, \texttt{action}, \texttt{before}, \texttt{after} và \texttt{createdAt}.

\subsubsection{Thông tin actor và request}

Actor và request ID được xác định ở backend, không lấy trực tiếp từ body request, để tránh giả mạo lịch sử.

\subsubsection{Snapshot trước và sau thay đổi}

Snapshot \texttt{before}/\texttt{after} dùng \texttt{JSONB}, đủ để so sánh thay đổi nhưng không lưu bí mật hoặc dữ liệu cá nhân không cần thiết.

\subsubsection{Tính chất append-only}

Ứng dụng chỉ thêm audit entry mới, không cung cấp luồng sửa hoặc xóa bản ghi đã tạo.

\subsubsection{Lịch sử cấu hình dựa trên audit log}

History API dựng lại diễn biến cấu hình từ các audit entry liên quan thay vì duy trì bảng version riêng.

\subsection{Ràng buộc toàn vẹn dữ liệu}

Ràng buộc dữ liệu được đặt ở cả Prisma/database và tầng validation để giữ cấu hình hợp lệ trước khi đi vào engine.

\subsubsection{Ràng buộc khóa và tính duy nhất}

\begin{table}[H]
    \centering
    \caption{Các ràng buộc duy nhất tiêu biểu}
    \label{tab:unique-constraints}
    \begin{tabular}{p{0.30\textwidth} p{0.55\textwidth}}
        \toprule
        \textbf{Thực thể} & \textbf{Ràng buộc} \\
        \midrule
        \texttt{Project} & \texttt{key} duy nhất toàn hệ thống. \\
        \texttt{Environment} & \texttt{projectId, key} duy nhất trong project. \\
        \texttt{FeatureFlag} & \texttt{projectId, key} duy nhất trong project. \\
        \texttt{FlagEnvironmentConfig} & \texttt{featureFlagId, environmentId} duy nhất. \\
        \texttt{FlagGroupConfig} & \texttt{groupId, environmentId} duy nhất. \\
        \bottomrule
    \end{tabular}
\end{table}

\subsubsection{Ràng buộc quan hệ tham chiếu}

Khóa ngoại bảo đảm flag, config, rule, group và audit entry không trỏ đến project hoặc environment không tồn tại.

\subsubsection{Ràng buộc giá trị phần trăm rollout}

Percentage rollout phải nằm trong miền hợp lệ, thường từ 0 đến 100.

\subsubsection{Ràng buộc không sửa hoặc xóa audit log}

Audit log không có API update/delete; triển khai production có thể bổ sung trigger database để cưỡng chế thêm.

\subsubsection{Chiến lược archive thay cho xóa vật lý}

Flag không còn dùng được archive thay vì xóa vật lý để giữ lịch sử và tránh phá vỡ audit/history.
